\chapter{Preface}
\label{ch:preface}

I wanted a book in which a graph that composes successfully could still be
the beginning of the problem. The interesting questions come after the
services start: why one query issues too many database statements, why the
right records arrive in the wrong positions, why guarding a list leaves
another route open, and why one failed field can erase data from a healthy
service. A diagram of several services behind one endpoint does not answer
those questions. A request, its response and the work behind it can.

This book builds a conference API in C\# with Hot Chocolate, then splits it
behind a Cosmo router. The example stays small enough that every row has a
name and every extra statement can be accounted for. Sessions, speakers and
ratings give it two different kinds of join. Search later supplies a page
whose ordering depends on data the schedule service does not hold. The
point is to make the boundaries observable, not to present a production
conference system.

I assume you have shipped an ASP.NET Core service, can read C\# and JSON,
and have seen a GraphQL schema and query. You do not need federation
experience. The first chapters establish one working service before adding
a composer or router. If several teams are not blocked on changing the same
API, that service may be all you need. Chapter~\ref{ch:do-you-need-this}
makes the adoption decision before the book asks you to install the tools.

Read the build-along chapters in order. A later listing assumes the state
left by the earlier one, including its seed data and generated loaders.
Temporary failures are exercises with a return path: restore the working
state before moving on. For a lookup, use the appendices. They collect the
toolchain, compatibility differences, directive vocabulary, composition
diagnostics and complete source through chapter~\ref{ch:ownership}, so a
reader does not have to reconstruct the final files from successive edits.

The measurements are deliberately narrow. A statement count describes the
query and data printed beside it; it is not a latency benchmark or a promise
about your database. An HTTP request count describes the boundary where it
was observed. Where the specification, library and router disagree, I name
the disagreement and keep the observed result. A passing composition is a
useful check, not proof that every client operation can be served.

The package versions matter here because the book prints APIs and diagnostic
text, not just concepts. Appendix~\ref{app:toolchain} gives the reproducible
baseline. The Hot Chocolate 14 callouts are for readers maintaining older
code; that release no longer receives security fixes. They are not a
recommendation to begin a new service on it. Nor are this book's local token,
SQLite reset or static router configuration production deployment advice.

This manuscript currently develops the example through
chapter~\ref{ch:ownership}. The later testing and operational chapters listed
in the contents are still to come. The source appendix describes the working
state reached here, not an implementation of those unwritten chapters.

Start with the field your team cannot ship. If federation removes that
release bottleneck, the rest of the book gives you a small graph whose
failures you can reproduce before they belong to a much larger one.
